\begin{frame}{Salud y Colpensiones elevarían el gasto en 1,9 puntos del PIB}
    \begin{itemize}
        \item Se puede lograr un recorte de más de dos puntos del PIB revirtiendo los aumentos en la inversión, en el subsidio a los combustibles y en otras transferencias. Sin embargo, persistirán presiones adicionales:
    \end{itemize}

    \begin{table}[H]
    \centering
    \small
    \setlength{\tabcolsep}{5pt}
    \begin{tabular}{l l c r r r}
    \hline
    \textbf{Rubro} & \textbf{Fuente} & \textbf{Período} & \textbf{Inicio} & \textbf{Final} & \textbf{Cambio} \\
    \hline
    Salud & MFMP 2026 & 2026--36 & 3,4 & 4,5 & \textbf{+1,1} \\
    Colpensiones & MFMP 2026 & 2026--36 & 1,7 & 2,5 & \textbf{+0,8} \\
    \textbf{Total} & \textbf{MFMP 2026} & \textbf{2026--36} & \textbf{5,1} & \textbf{7,0} & \textbf{+1,9} \\
    \hline
    \end{tabular}
    \\[0.35em]
    \scriptsize Niveles y cambios en puntos del PIB. Los aumentos pueden sumarse porque corresponden a rubros distintos, con la misma fuente y el mismo período.\\[0.2em]
    \tiny Fuente: MFMP 2026, gráfs. A.1.5--A.1.6.
    \end{table}
    \note{[Sources] MFMP 2026, Apéndice 1, gráficos A.1.5 y A.1.6: https://www.minhacienda.gov.co/documents/d/portal/mfmp-2026?download=true. Valores 2026--2036: aportes a salud de 3,4\% a 4,5\% del PIB (+1,1 pp); aportes a Colpensiones de 1,7\% a 2,5\% (+0,8 pp). La suma pasa de 5,1\% a 7,0\% del PIB (+1,9 pp). Consultado el 20 de agosto de 2026.}
\end{frame}

\begin{frame}{Contener el gasto en salud y pensiones}
    \begin{itemize}
        \item La consolidación fiscal difícilmente será posible sin contener las presiones de gasto en salud y pensiones.
        \item Para \textbf{salud} se requiere:
        \begin{itemize}
            \item Buscar fuentes de financiación por fuera del PGN.
            \item Retornar a un plan de beneficios explícitos.
        \end{itemize}
        \item Para \textbf{pensiones} se requiere:
        \begin{itemize}
            \item Reformas paramétricas.
            \item Gravar las pensiones con el impuesto de renta.
            \item Revisar el mandato constitucional de pensión mínima de un salario mínimo.
        \end{itemize}
    \end{itemize}
\end{frame}

\begin{frame}{El aumento del SGP no es necesariamente gasto adicional}
    \begin{itemize}
        \item El CARF proyecta que el SGP pase de 4,4\% del PIB en 2026 a 5,8\% en 2037: un aumento bruto de 1,4 puntos.
        \item El MFMP supone neutralidad fiscal porque la inversión disminuye en la misma magnitud. El CARF no puede verificar ese supuesto mientras no se definan las competencias por trasladar y su costo.
        \item Parte del SGP financia salud. Si los recursos adicionales sustituyen aportes de la Nación para las mismas obligaciones, sumar el aumento del SGP al de salud contaría dos veces parte del gasto.
    \end{itemize}
    \textbf{La presión neta dependerá de cuánto gasto del GNC sea efectivamente sustituido.}\\[0.25em]
    \scriptsize Fuentes: CARF (2026), pp. 26--27; Ministerio de Salud.
    \note{[Sources] CARF, Documento técnico sobre el MFMP 2026, pp. 26--27 y gráfica 8: el escenario inercial proyecta que el SGP pase de 4,4\% del PIB en 2026 a 5,8\% en 2037 (+1,4 pp). El CARF incorpora el Acto Legislativo 03 de 2024 desde 2028, pero no los efectos de la Ley de Competencias porque el articulado no permite identificar las funciones por trasladar ni su costo; el MFMP supone neutralidad fiscal mediante una reducción de la inversión equivalente al aumento del SGP: https://www.carf.gov.co/documentos-tecnicos/informes-y-documentos-tecnicos. Ministerio de Salud: el SGP es una fuente de financiación de salud; por ello, un traslado fiscalmente neutral debe sustituir gasto nacional asociado a las mismas obligaciones y no adicionarse a él: https://minsalud.gov.co/Proteccion-Social/financiamiento/Paginas/financiacion-aseguramiento.aspx y https://www.minsalud.gov.co/proteccionsocial/Financiamiento/Paginas/sgp-salud-publica.aspx. Consultados el 20 de agosto de 2026.}
\end{frame}

\begin{frame}
    \begin{itemize}
        \item Inevitablemente, una reducción de al menos 3 puntos del PIB en el déficit primario deberá realizarse a través de aumentos en los ingresos.
        \item No obstante, \textbf{la fórmula del SGP} implica que cada 100 pesos de ingresos adicionales generan eventualmente al menos 25 pesos de gasto.
        \item Por lo tanto una reducción del déficit primario de \textbf{3 puntos del PIB} a través de un aumento de los ingresos en realidad requiere un \textbf{recaudo} adicional de  \textbf{4 puntos del PIB}.
        \item Si la Ley de Competencias genera gasto adicional no compensado, el aumento del recaudo tendría que ser aún mayor.
    \end{itemize}
\end{frame}

\begin{frame}{Resumiendo}
\begin{itemize}
    \item La \textbf{consolidación fiscal} requiere:
    \begin{itemize}
        \item Recortes en la inversión y otras transferencias.
        \item Contención del gasto en pensiones y salud.
        \item Posponer la Ley de Competencias.
        \item Reformular al SGP como porcentaje del PIB.
        \item Marchitar beneficios tributarios.
        \item Reformar administrativamente a la DIAN.
%        \item Implementar un ajuste gradual.
    \end{itemize}
    \item La gradualidad requerirá de una \textbf{nueva regla fiscal}.
\end{itemize}    
\end{frame}

\begin{comment}
\begin{frame}{¿Por qué siguen prestando los inversionistas?}
\begin{itemize}
    \item Los inversionistas \textbf{no son ingenuos}: es una combinación de restricciones institucionales e incentivos perversos.
    \begin{itemize}
        \item \textbf{AFP y bancos domésticos}: atrapados por regulación — no tienen alternativa doméstica de escala comparable a los TES.
        \item \textbf{Inversionistas extranjeros}: apuestan al \textit{carry} mientras no haya un \textit{trigger} claro de default, esperando salir antes de que el problema explote.
        \item \textbf{Represión financiera progresiva}: muchos no anticipan un default duro sino tasas reales negativas y licuación gradual de la deuda.
    \end{itemize}
    \item La paradoja del \textbf{Banrep independiente}: si no monetiza, el ajuste o el default son las únicas salidas. Si cede, se destruye el ancla que hace atractivos los TES.
    \item \textbf{Conclusión}: la presión para ajustar no vendrá de adentro — vendrá de un choque externo (caída del petróleo, alza de tasas globales, rebaja a \textit{high yield}).
\end{itemize}
\end{frame}

\begin{frame}{Lecciones de casos comparados}
\begin{itemize}
    \item En ningún caso el ajuste vino de adentro, sin presión externa. Lo que varía es si el país llega a ese momento con instituciones suficientemente fuertes para \textbf{negociar} el ajuste en vez de \textbf{sufrirlo}.
\end{itemize}
\begin{table}[H]
\centering
\small
\begin{tabular}{l p{2.8cm} p{2.2cm} p{2.2cm}}
\hline
\textbf{Caso} & \textbf{Mecanismo} & \textbf{\textit{Trigger}} & \textbf{Costo} \\
\hline
Ecuador 2000 & Dolarización forzada + default & Crisis bancaria + choques externos & Devastador \\
Ecuador 2020 & Default + renegociación FMI & COVID-19 & Moderado \\
Costa Rica 2018--24 & Reforma fiscal negociada & Rebaja a \textit{high yield} & Alto pero manejable \\
\textbf{Colombia (?)} & \textbf{?} & \textbf{Por determinar} & \textbf{Por determinar} \\
\hline
\end{tabular}
\\[0.3em]
\footnotesize Fuente: elaboración propia con base en FMI, OCDE y BFI.
\end{table}
\begin{itemize}
    \item \textbf{Ecuador}: el ajuste políticamente imposible se volvió inevitable — el \textit{default} de 2008 fue de \textit{unwilling to pay}, no de incapacidad fiscal.
    \item \textbf{Costa Rica}: el ajuste pasó en parte porque nadie entendió completamente lo que se aprobaba en 2018 — y aun así requirió una rebaja a \textit{high yield} como catalizador.
\end{itemize}
\end{frame}
\end{comment}
